% Section 11 -- drop-in replacement for "Provide the conclusion here".
% Plain names (DUTS, Greedy-Swap, Preference-Swap, Santos-Small, ...); substitute your own
% macros if the source defines them. Note: p.11-12 write "Santo-Small" in the prose and
% Table 2/3 -- fix to "Santos-Small" throughout, or keep whichever you intend.

\section{Conclusion and Future Work}

We introduced \emph{Distribution-Aware Unionable Table Search} (DUTS): given a query table
and a distributional requirement over one of its categorical attributes, retrieve $k$
unionable tables whose multiset union with the query satisfies that requirement while
maximizing unionability. Because the requirement constrains the \emph{union} rather than any
individual table, the problem is aggregation-sensitive, and it does not decompose into
independent per-table tests. We proved that both its decision and optimization versions are
NP-hard by reduction from cardinality-constrained knapsack.

We studied two paradigms for solving it. The \emph{unionability-first} paradigm retrieves a
high-unionability result and then repairs it by swapping tables; we showed that computing the
cheapest such repair is itself NP-hard by reduction from Subset Sum, and developed a greedy
swap heuristic together with a preference-based filter that provably never discards an
optimal repair. The \emph{distribution-first} paradigm, our two-stage solution, pushes the
requirement into retrieval instead: Stage 1 selects a distribution-optimal pool of $\alpha k$
candidates by solving a $0$--$1$ linear-fractional program exactly with Dinkelbach's
parametric method, and Stage 2 selects the final $k$ tables by maximizing unionability over
that pool as a linear $0$--$1$ integer program, since fixing the threshold $\tau$ linearizes
the ratio constraint. Stage 2 therefore never scores more than $\alpha k$ candidates,
independent of the size of the data lake.

Our evaluation shows that the distribution-first paradigm is the stronger of the two. DUTS
returns a distributionally feasible result for the large majority of queries (39/48, 84/92,
and 122/142 on Santos-Small, TUS-Small, and TUS-Large) while retaining most of the
distribution-oblivious baseline's precision and aggregate unionability; the swap heuristics
either satisfy far fewer queries or pay a substantially larger quality cost to satisfy them.
The two paradigms also fail for different reasons. Every DUTS failure we observed is
attributable to retrieval returning a pool that contains no feasible set, which is what its
formulation predicts: given a pool that does contain one, the optimization finds it. The
heuristics, in contrast, fail on queries whose retrieved pool does contain a feasible
solution, as the greedy order or the filtering step can discard the tables such a solution
needs.

The optimization machinery is inexpensive at scale. Dinkelbach's method converges in fewer
than three iterations on average on every benchmark and pool size we measured, including
candidate sets of roughly 5{,}600 tables, and Stage 1 and Stage 2 together stay in the
single-digit milliseconds. End-to-end latency remains at roughly 400\,ms or below on a
one-million-table lake, and is dominated not by the optimization but by the inverted-index
probe, whose cost is output-sensitive in the length of the posting list it must materialize.

\paragraph{Future work.}
Our results point to retrieval, rather than optimization, as the component that limits both
quality and latency, and this shapes the directions we consider most promising.

\emph{Output-sensitive retrieval.} At the one-million-table tier the overlap probe accounts
for roughly 84\% of retrieval time, because it enumerates every $(\text{table},
\text{column})$ pair holding a distribution value even though only a small fraction survives
intersection with the semantic candidates. Compressed posting lists, bitmap intersection, and
early termination would let the probe pay for what survives rather than for what it reads.

\emph{Constraints in candidate generation.} A more fundamental direction is to make the index
itself yield distribution-feasible candidates, rather than filtering for them afterwards.
Recent work embeds diversity constraints into graph traversal for approximate nearest-neighbor
search, but that setting fixes the constraint before indexing. Ours does not: the distribution
attribute $V_D$ and the value set $\mathcal{M}$ are query-time parameters, and which values a
user treats as distributionally salient is not known when the index is built. Designing an
index that supports constraints chosen at query time is an open problem.

\emph{Richer requirements.} Our formulation targets a single proportion $F^\star$ for one value
set over one attribute. Requiring a full target histogram over $Dom(V_D)$, or requirements over
several attributes at once (as intersectional notions of representation demand), changes the
structure of the constraint and would need different machinery in both stages.

\emph{The Stage 1 objective.} Stage 1 currently maximizes the achieved proportion without a
bound, which fills the pool with the most distributionally extreme tables. A satisficing
variant that stops once $\tau$ is reachable and spends the remaining freedom on a secondary
criterion, such as pool size or unionability potential, may hand Stage 2 a better pool; how to
choose $\alpha$ adaptively is a related question.

\emph{Cost-aware selection.} When tables are acquired from a data market at a price, the
natural problem is to satisfy the distributional requirement within a budget. Adapting
tuple-level distribution-tailoring frameworks to table-level selection, with unionability loss
in the role of acquisition cost, would give both a new baseline and a practical variant of the
problem.
